% Author: JoonHo Lee (jlee296@ua.edu)
\begin{table}[htbp]\centering
\caption{Information indices as approximations to fitted EAP and WLE reliability}
\label{tab:s4-correspondence}
{\small\begin{tabular}{llrrrr}
\toprule
Score & Comparison & Favours second & Median advantage & Sign $p$ & Status \\
\midrule
EAP & $\wbar$ vs $\rhopsd$ & 28/48 & $6.18\times 10^{-4}$ & $0.3123$ & unresolved \\
EAP & $\rhotilde$ vs $\rhopsd$ & 37/48 & $0.001779$ & $2.22\times 10^{-4}$ & supported \\
WLE & $\rhopsd$ vs $\wbar$ & 43/48 & $0.001452$ & $<10^{-4}$ & supported \\
\bottomrule
\end{tabular}}
\floatnote{The 48 cells cross Rasch/2PL, two ability shapes, two lengths, two person counts and three targets, with 200 response replications each. Counts favoring the second index, median absolute-distance advantages and sign tests summarize those cells. The EAP comparison between pointwise and inverse information is inconclusive ($p=.3123$). These are comparisons with fitted TAM coefficients, not identities for population reliability.}
\end{table}
